% !PRINTED 706
\weberpage{706}
Da diese Funktion hiernach keine negativen Potenzen von $z$ enthalten kann, so ergibt sich aus (13) der zu beweisende Satz:
\[
\overset{i}{\Sigma} a^{(i)}_{-1} = 0,
\]
der mit der besonderen Annahme über die Variable $z$ nichts mehr zu tun hat.

Als spezielle Anwendung folgt hieraus, daß ein Differential zweiter Gattung kein Residuum hat, und daß die beiden Residuen der Differentiale dritter Gattung gleich und entgegengesetzt sind.

\gesperrt{Ein eigentliches Differential} $d\sigma$, wie wir schon gesehen haben, \gesperrt{hat kein Residuum}, und die Residuen des \glqq\gesperrt{logarithmischen Differentials}\grqq{} $d\sigma:\sigma$ sind ganze Zahlen, nämlich die Ordnungszahlen in $z$.

Wir wollen diese Betrachtungen mit dem Beweis des folgenden Satzes beschließen:

4. \gesperrt{Nennen wir Differentiale zweiter Gattung linear abhängig, wenn eine lineare Verbindung von ihnen und von Differentialen erster Gattung mit konstanten Koeffizienten einem eigentlichen Differential gleich ist, so sind höchstens $p+1$ Differentiale zweiter Gattung immer linear abhängig.}

Um dies zu beweisen, nehmen wir ein $p$-Eck zweiter Gattung:
\weq{18}{\mathfrak P_1 \mathfrak P_2 \ldots \mathfrak P_p}
mit $p$ verschiedenen Punkten. Ist dann $\mathfrak P$ ein gegebener Punkt, so setzen wir:
\weq{19}{\mathfrak U = \mathfrak P \mathfrak P_1 \mathfrak P_2 \ldots \mathfrak P_r}
% [sic] printed P_r (Fraktur) as the last factor of (19), a misprint for P_p: the p-gon (18) is P_1 P_2 ... P_p
und erhalten eine Klasse $B$, in deren Teiler jedenfalls $\mathfrak P$ nicht aufgeht, weil sonst $\mathfrak P_1 \mathfrak P_2 \ldots \mathfrak P_n$
% [sic] printed P_1 P_2 ... P_n (Fraktur), a misprint for P_1 P_2 ... P_p: the polygon meant is the p-gon (18)
nach $\S\,199$, 7. von der ersten Gattung sein müßte. Wir können also eine Funktion in $\Omega$:
\weq{20}{z = \frac{\mathfrak A}{\mathfrak U}}
bestimmen, in deren Untereck jedenfalls der Punkt $\mathfrak P$ enthalten ist. Wir bilden nun das eigentliche Differential
\weq{21}{dz = \frac{\mathfrak Z}{\mathfrak U^2},}
in dessen Zähler $\mathfrak Z$ (dem Verzweigungspolygon nach $z$) der Punkt $\mathfrak P$ nicht aufgeht, weil er in $\mathfrak U$ nur einfach aufgeht, und wenn wir $dz$ nach $\S\,200$, 4. durch Differentiale erster, zweiter und dritter
%CONT>>
